\documentclass[exercice,a4]{thedocument}%
\usepackage[french]{babel}%
\usepackage{pstricks}%
\startdatakeys
    \source{22GENMATNC1}
    \BO{2026}
    \cycle{4}
    \competencesDuSocle{CA,RA}
    \objectifsApprentissage{C4C2bc}
    \motsCles{probabilité}
\enddatakeys
\begin{document}%

\medskip

\begin{minipage}{0.6\linewidth}
Gabriel lance deux fois de suite un dé équilibré à quatre faces numérotées de 1 à 4 et il relève le numéro qui figure sur la face cachée du dé.

Si Gabriel obtient 2 au premier lancer puis 4 au second, il note (2~;~4).
\end{minipage}\hfill
\begin{minipage}{0.35\linewidth}
\begin{center}
\psset{unit=1cm}
\begin{pspicture}(4.5,3.7)
\pspolygon[fillstyle=solid,fillcolor=red](0,0.2)(1.95,3.6)(4.15,1.1)(4,0)
\psline[linewidth=2pt,linecolor=white](1.95,3.6)(4,0)
\rput(2,1.5){\white \Huge 3}
\rput{28}(3.4,1.5){\white \Huge 1}
\end{pspicture}
\end{center}
\end{minipage}

\medskip

\begin{enumerate}
\item Gabriel a noté (3 ; 2).
	\begin{enumerate}
		\item Quel numéro a-t-il obtenu au premier lancer ?
		\item Quel numéro a-t-il obtenu au second lancer ?
	\end{enumerate}
\item Quelles sont les $16$ issues possibles de ce jeu ?
\item Que dire de l'évènement $A$ : \og Obtenir 1 en additionnant les deux numéros obtenus\fg{} ?

L'évènement $B$ : \og Obtenir 7 en additionnant les deux numéros obtenus \fg{} peut être réalisé avec l'issue (3~;~4) ou avec l'issue (4~;~3).
\item Donner les quatre issues possibles qui réalisent l'évènement $C$ : \og Obtenir 5 en additionnant les deux numéros obtenus \fg.
\item Quelle est la probabilité que l'évènement $C$ se réalise?
\end{enumerate}

\bigskip

\showthesource
\end{document}
